\documentclass[12pt,a4paper]{article}
\usepackage[utf8]{inputenc}
\usepackage{geometry}
\geometry{a4paper, margin=1in}
\usepackage{hyperref}
\usepackage{graphicx}
\usepackage{setspace}
\usepackage{cite}
\usepackage[most]{tcolorbox}

\title{\textbf{Unlocking the Microscale: How Digital Rock Physics Can Secure Supply Chains for UK and European Critical Minerals}}
\author{Hannah P. Menke \and Maja Rücker}
\date{\today}

\begin{document}

\maketitle

\begin{abstract}
The United Kingdom and Europe face acute supply chain vulnerabilities for Critical Raw Materials (CRMs) essential to the energy transition and defense. While policy frameworks such as the UK Critical Minerals Strategy and the EU Critical Raw Materials Act (CRMA) set aggressive targets for domestic extraction, processing, and recycling, the technical means to extract maximum value from increasingly complex, low-grade ores still lag behind. The ``Digital Rocks'' or Digital Rock Physics (DRP) community—historically anchored in oil and gas—already offers deep expertise in multi-scale, multi-modal imaging (including CT, nano-CT, MRI, PET, and FIB-SEM) and pore-scale flow simulation. This white paper argues for a coordinated pivot of that expertise toward ore characterization, extraction technologies such as Direct Lithium Extraction (DLE), and mine waste valorization. Rather than reviewing critical minerals policy or geometallurgy in general, it proposes a DRP-led framework linking microstructural imaging, geometallurgy, recycling, and porous-material design. Used this way, DRP could shorten evaluation times, reduce environmental impacts, and strengthen the UK's and Europe's technological position in the race for secure, sustainable critical minerals.
\end{abstract}

\section{Context: The Geopolitical Imperative and the Global Race}
The global transition to clean energy technologies has triggered rapidly rising demand for critical minerals such as lithium, cobalt, rare earth elements (REEs), and graphite. Recent review literature places this challenge in a wider geopolitical and technological context, emphasizing both resource security and processing capability as core bottlenecks \cite{reich_critical_2025}. Global lithium demand, for example, surged by nearly 30\% in 2024, even as mining exploration activity plateaued and investment growth slowed markedly \cite{iea_outlook_2025}. Consequently, achieving higher recovery and better resource efficiency from existing deposits is increasingly important.

Compounding this issue is the hyper-concentration of the market. The IEA reports that the average market share of the top three refining nations for key energy minerals rose to 86\% in 2024 \cite{iea_outlook_2025}. In response, the EU Critical Raw Materials Act (CRMA) established 2030 benchmarks requiring 10\% extraction, 40\% processing, and 25\% recycling to occur domestically \cite{eu_crma_2024}. Concurrently, the UK's strategy focuses on resilient supply chains, domestic capability, and leveraging world-leading academic strengths \cite{uk_strategy_2022}.

Globally, the United States is also funding advanced mining innovation through programs such as ARPA-E's MINER portfolio and the DOE's \$80M ``Mine of the Future'' Proving Ground Initiative \cite{doe_arpae, doe_mine_future_2025}. The UK and Europe should mobilize their existing DRP communities if they want to remain competitive while also building productive transatlantic collaborations. As ore grades fall, the ``easy'' deposits have largely vanished, and the new frontier of mining requires far greater micro-scale precision \cite{calvo_copper_2016}.

\section{The Digital Rocks Paradigm}
Digital Rock Physics (DRP) uses a suite of 3D and 4D imaging technologies to map the structure of porous materials. While often associated primarily with X-ray micro-computed tomography ($\mu$CT), DRP spans a much broader set of techniques, including medical and industrial CT, nano-CT, Positron Emission Tomography (PET), Magnetic Resonance Imaging (MRI), Nuclear Magnetic Resonance (NMR), Focused Ion Beam Scanning Electron Microscopy (FIB-SEM), and Helium Ion Microscopy (HIM). Coupled with image processing, AI-driven segmentation, and numerical simulation, these methods form a unified platform for tracking mass, charge, and fluid flow \cite{drp_blunt_2013}.

Historically, this toolkit was developed to optimize hydrocarbon extraction. However, as the energy transition accelerates, there is a strong case for pivoting this existing infrastructure—including synchrotrons, lab-scale scanners, and supercomputing clusters—towards hard-rock mining and mineralogy. 

This argument sits adjacent to, but is distinct from, recent reviews of integrated geometallurgy and applied mineralogy across the mine life cycle \cite{lishchuk_integrated_2020, becker_applied_2023}. Those literatures show why mineralogical heterogeneity, texture, and sustainability matter. The contribution here is to argue that DRP offers a practical technical bridge across those domains when imaging, flow simulation, and data infrastructure are considered together.

Crucially, the UK and Europe already host much of the infrastructure needed for this shift. Diamond Light Source supports tomography and correlative imaging across beamlines including I13, I12, I14, and DIAD; ESRF's BM05 imaging line provides high-energy XCT and phase-contrast imaging for 3D materials characterization; and PSI's TOMCAT platform spans fast micrometer-scale to sub-100-nm tomography for geoscience and energy-materials research \cite{diamond_tomography_2026, esrf_bm05_2026, psi_tomcat_2025}. The strategic challenge is therefore less about building capability from scratch and more about coordination, shared data standards, and translation into mineral-processing workflows.

\begin{figure}[htbp]
\centering
\small
\setlength{\tabcolsep}{4pt}
\renewcommand{\arraystretch}{1.15}
\begin{tabular}{|p{0.225\textwidth}|p{0.225\textwidth}|p{0.225\textwidth}|p{0.225\textwidth}|}
\hline
\parbox[t]{0.215\textwidth}{\textbf{Exploration}\newline 3D texture priors\newline AI training data\newline rapid screening of complex ores} &
\parbox[t]{0.215\textwidth}{\textbf{Characterization}\newline mineral associations\newline liberation metrics\newline correlative chemistry + structure} &
\parbox[t]{0.215\textwidth}{\textbf{Processing \& Extraction}\newline breakage behavior\newline leach-flow pathways\newline porous sorbent design} &
\parbox[t]{0.215\textwidth}{\textbf{Recycling \& Tailings}\newline sensor calibration\newline black-mass texture\newline secondary resource mapping} \\
\hline
\end{tabular}
\caption{Where Digital Rock Physics can add decision-relevant information across the critical-minerals value chain.}
\label{fig:drp_value_chain}
\end{figure}

Taking inspiration from the US DOE National Energy Technology Laboratory’s ``Digital Rocks Portal,'' there is a strong case for a federated European ``Digital Ore Database'' that shares multi-modal tomographic data on critical-mineral deposits and broadens access to high-resolution structural data \cite{netl_digitalrocks, digitalrocks_portal_2023}.

For the purposes of this paper, it is useful to distinguish between broad strategic categories of critical minerals and a limited set of representative examples. This framing retains a policy-level perspective while illustrating where DRP may have the greatest practical relevance in UK and European settings.

\begin{table}[htbp]
\centering
\small
\setlength{\tabcolsep}{4pt}
\renewcommand{\arraystretch}{1.15}
\begin{tabular}{|p{0.18\textwidth}|p{0.18\textwidth}|p{0.31\textwidth}|p{0.25\textwidth}|}
\hline
\textbf{Category} & \textbf{Example minerals} & \textbf{Why DRP helps} & \textbf{Why relevant in UK/Europe} \\
\hline
Battery materials & Lithium, cobalt, nickel, manganese, graphite & 3D texture mapping, liberation analysis, pore-scale transport, black-mass characterization, porous sorbent design for DLE & Supply-chain resilience for batteries and electrification, plus growing recycling and midstream-processing ambitions \\
\hline
Magnet and high-tech materials & Rare earth elements, gallium, germanium & Correlative imaging for finely disseminated phases, automated mineralogy, selective separation workflow design & Strategic importance for magnets, electronics, defense, and limited tolerance for supply disruption \\
\hline
Technology and alloy metals & Tungsten, tin, antimony & Ore-texture characterization, selective breakage studies, tailings re-evaluation, sensor calibration for sorting & Relevant to legacy European mining districts and to reprocessing of historical mine wastes \\
\hline
Secondary resources & Tailings, battery black mass, industrial residues & Multi-scale characterization of heterogeneous feedstocks, recovery-pathway design, calibration of sorting and recycling processes & Particularly compatible with UK/EU circular-economy goals, existing industrial residues, and recycling policy frameworks \\
\hline
\end{tabular}
\caption{Representative mineral categories and examples that show where DRP can add value without reducing the paper to a commodity-by-commodity inventory.}
\label{tab:mineral_categories}
\end{table}

\section{Pillar I: Exploration and Advanced Characterization}
The implementation of DRP can shift mineral exploration strategies, particularly by supporting AI-based exploration. Recent review literature identifies AI-enabled geological exploration as a rapidly developing frontier, with mineral-prospectivity mapping remaining one of its best-established use cases \cite{yang_ai_exploration_2024, ai_mineral_prospectivity_2011}. DRP provides ultra-high-resolution, 3D ground-truth datasets that can improve these predictive models. It can also shorten the path from discovery to commercial evaluation by enabling faster orebody assessment using advanced sensing and analytical methods \cite{doe_arpae}.

Moving beyond traditional 2D thin sections, multi-modal tomographic approaches from macro-CT to nano-CT can quantify 3D mineral associations, grain sizes, and liberation characteristics. When paired with SEM-EDX-based automated mineralogy, these correlative workflows can locate finely disseminated or texturally complex critical minerals that might otherwise escape detection \cite{buyse_correlative_2023}. Emerging 3D liberation models further suggest a route from ore texture to process-relevant particle behavior \cite{guntoro_liberation_2021}. Together, these data-rich models support predictive geometallurgy by integrating sparse process data into spatial orebody models before full-scale mining, improving flowsheet design from the outset \cite{lishchuk_predictive_2019}.

\begin{tcolorbox}[enhanced, colback=blue!3!white, colframe=blue!35!black, title=\textbf{Illustrative Case Study: Correlative Imaging for Complex Pegmatites}, fonttitle=\bfseries, arc=1.5mm, boxrule=0.8pt]
Buyse et al. combined X-ray micro-CT with SEM-EDX automated mineralogy to resolve internal mineral relationships in Sn-Nb-Ta pegmatites and build a mineral library for 3D interpretation \cite{buyse_correlative_2023}. The implication for critical minerals is straightforward: similar workflows can be adapted to lithium-bearing pegmatites and other texturally complex ores where 2D sections alone are an incomplete guide to processing behavior.
\end{tcolorbox}

\section{Pillar II: Optimizing Mineral Processing and Extraction}
The extraction of critical minerals is notoriously energy- and water-intensive, yet DRP offers several pathways to optimize these core processes. Comminution (crushing and grinding), for example, is exceptionally energy-intensive, accounting for large amounts of global electricity consumption and often exceeding 50\% of a single mine site's power demand \cite{norgate_comminution_2010}. DRP enables microstructurally informed mineral liberation and breakage studies that can support more selective comminution strategies \cite{kyle_processing_2015}. 

Similarly, in operations such as in-situ and heap leaching, engineers can use dynamic imaging (e.g., PET or MRI) and pore-scale flow simulation to map the transport of lixivium \cite{mri_heap_leach_2015}. That kind of process understanding can improve permeability management and dissolution efficiency while reducing unnecessary excavation. DRP is also relevant to emerging direct extraction technologies. Engineered porous materials such as metal-organic frameworks (MOFs), ion sieves, and ion-selective membranes are increasingly being explored for lower-footprint separations, including lithium extraction from brines \cite{mofs_lithium_2026, mojid_dle_2024, baird_dle_2024}. In these systems, the key imaging target is often not the brine itself but the pore architecture, accessibility, degradation, and cycling behavior of the sorbents or membranes that control selectivity and throughput.

\section{Pillar III: Tailings Valorization and the Circular Economy}
Developing a circular economy is a cornerstone of both UK and EU critical minerals strategies. Legacy mine tailings frequently contain significant, unrecovered amounts of critical minerals such as tellurium, antimony, and tungsten. Recent reviews highlight both the strategic importance of reprocessing mine tailings and the opportunity to recover critical materials from them \cite{sarker_tailings_2022, araya_tailings_2020, nwaila_minewaste_2021}. 

To recycle these materials effectively, advanced sorting mechanisms need to be deployed and calibrated. Sensor-based sorting reviews in mineral processing point to preconcentration, waste reduction, and multi-sensor fusion as key opportunities for more selective treatment of heterogeneous feedstocks \cite{peukert_sorting_2022}. High-resolution 3D imaging can serve as a micro-scale calibration tool for these industrial sensors. Multi-scale DRP methods can also be applied to unconventional secondary feedstocks such as coal refuse, fly ash, and acid mine drainage. In addition, techniques such as FIB-SEM and nano-CT can characterize e-waste and battery ``black mass'' at sub-micron scales to optimize both mechanical and chemical recycling pathways by mapping structural degradation and elemental composition \cite{battery_blackmass_2021}.

\begin{tcolorbox}[enhanced, colback=green!3!white, colframe=green!35!black, title=\textbf{Illustrative Case Study: Black-Mass Characterization}, fonttitle=\bfseries, arc=1.5mm, boxrule=0.8pt]
Vanderbruggen et al. showed that automated mineralogy can characterize the composition and texture of spent lithium-ion batteries in ways that are directly relevant to recycling process design \cite{battery_blackmass_2021}. For a DRP-led critical-minerals agenda, this matters because it moves ``urban mining'' from bulk assay toward texture-aware process optimization, where liberation, coating residues, and phase associations all influence recovery.
\end{tcolorbox}

\section{Next-Generation Digital Rocks: Expanding the Horizon}
As Digital Rock Physics matures, it should expand toward increasingly complex chemical and synthetic challenges. Foremost among these is the integration of chemistry with structure via correlative microscopy \cite{buyse_correlative_2023}. Digital rocks should evolve beyond physical characterization metrics such as porosity and permeability to provide explicit chemical context in 3D structures. The field should adopt multi-modal workflows that register 2D elemental data, automated mineralogy, and 3D structural volumes in a geometallurgical framework \cite{koch_geomet_2019, gu_automated_mineralogy_2014}. That chemical overlay is important for tracking elemental speciation and reactive transport during complex, multi-phase extraction processes.

The same shift also extends the DRP toolkit from natural geology to engineered materials. Applying these imaging and simulation techniques to synthetic porous media—such as MOFs, zeolites, ion-imprinted membranes, and battery electrodes—has growing relevance for both DLE and downstream clean-energy technologies. DRP techniques can be used to design these porous structures across length scales ranging from angstroms to centimeters, with the aim of regulating the flow of mass, charge, heat, pressure, and radiation \cite{farber_science_2025}. Tools such as NMR, used alongside macroscopic CT, can help optimize surface-area-to-volume efficiency in these synthetic domains.

\section{The AI Revolution in Digital Rocks: Surrogates and Upscaling}
Machine Learning (ML) is also reshaping Digital Rock Physics by addressing one of its longstanding bottlenecks: the computational cost of pore-scale simulations. Traditional DRP simulations—particularly for multi-phase reactive transport—often require prohibitive computing resources. Today, image-based deep-learning surrogates can predict key petrophysical properties directly from microstructural images \cite{graczyk_dl_porosity_2020}. More broadly, physics-informed ML is emerging as a practical route for accelerating coupled simulations while retaining physical consistency \cite{karniadakis_piml_2021}.

Furthermore, Physics-Informed Neural Networks (PINNs) can improve the reliability of these models. Unlike standard data-driven AI, PINNs natively embed physical laws (such as Navier-Stokes equations for fluid flow) directly into their loss functions. This mechanism helps ensure that the model predictions remain physically plausible even when tomographic data is sparse or noisy \cite{raissi_pinns_2019}.

ML also offers one credible route through the DRP ``upscaling'' challenge. Because high-resolution imaging such as nano-CT is confined to microscopic sample volumes, translating those insights to the macroscopic mine-site scale remains difficult. Modern ML workflows address this by identifying representative microscopic texture patterns across large datasets, running targeted physics simulations on those textures, and upscaling the resulting physical properties to the core and reservoir scale \cite{ml_upscaling_drp_2024}. In that sense, the integration of AI and physics has the potential to extend DRP from a microscopic analytical tool into a macro-scale predictive framework.

\section{Environmental Impact and Sustainability}
The deployment of DRP also carries significant implications for environmental sustainability. In the context of carbon storage (CCUS), DRP can support safer and faster study of $CO_2$ mineralisation in mafic and ultramafic rocks. Building on projects such as CarbFix in Iceland \cite{matter_carbfix_2009}, combining CCUS with extraction could secure domestic resources while contributing to permanent carbon sequestration and improving understanding of dynamic subsurface gas storage \cite{kelemen_peridotite_2020, menke_hydrogen_2024}. Fast X-ray imaging also improves our understanding of wettability and multiphase flow in porous media, insights that are transferable to flotation and leaching process design \cite{bartels_fast_xray_2017}.

\section{The Road Ahead: Challenges and Recommendations}
To fully realize the potential of DRP in securing critical minerals, several hurdles still need to be addressed. There is a clear need for standardization, particularly around imaging protocols and machine learning models across the European research and industrial community. Significant technical challenges also remain in upscaling, especially in bridging the gap from nano-scale structural data (FIB-SEM/nano-CT) to the core scale (Macro-CT/MRI/PET), and then to the mine-site scale.

To overcome these challenges, the recommendations below are framed as near-term actions that could be initiated within a single public funding cycle:

\begin{tcolorbox}[enhanced, colback=gray!3!white, colframe=blue!40!black, title=\Large\textbf{Key Policy Recommendations}, colbacktitle=blue!40!black, drop fuzzy shadow, fonttitle=\bfseries, breakable, arc=2mm, boxrule=1pt]
\textbf{1. Launch Translational Demonstrators by 2027:}\\
UKRI, Horizon Europe, and national innovation agencies should establish dedicated calls for DRP-to-processing demonstrators. Each project should require a tripartite consortium spanning imaging science, orebody owner or recycler, and process scale-up capability. A practical target would be three to five pilot projects that move from characterization to pilot-plant decisions within 24--36 months.

\vspace{2mm}
\textbf{2. Define a ``Digital Ore Passport'' Standard by 2027:}\\
Regulators, geological surveys, and research infrastructures should publish a minimum data specification for publicly supported mining, recycling, and tailings projects. At a minimum, this should include multi-modal imaging metadata, mineralogical ground truth, uncertainty estimates, and machine-readable process descriptors. The goal should be adoption across all publicly funded pilot projects by 2028.

\vspace{2mm}
\textbf{3. Stand Up a Federated European Digital Ore Database by 2028:}\\
The European Commission, national geological surveys, and major imaging facilities should create a shared repository with open metadata, benchmark datasets, and application programming interfaces for model development. A reasonable first milestone would be 100 well-annotated benchmark datasets covering hard-rock ores, tailings, and battery-recycling feedstocks within the first two years.

\vspace{2mm}
\textbf{4. Build Workforce Capacity with Cross-Disciplinary Training Centres:}\\
The UK and Europe should fund doctoral and industrial training centres focused on the intersection of imaging, machine learning, mineral processing, and recycling. The point is not only to train students, but to create a talent pipeline that can move between beamlines, laboratories, software teams, and pilot plants. A useful benchmark would be two to three centres and at least 40 jointly supervised researchers by 2030.
\end{tcolorbox}

\section{Conclusion}
Leveraging the existing Digital Rock Physics community gives the UK and Europe a plausible route to differentiation in the race for secure and sustainable critical minerals. The paper's contribution is not to review critical minerals, geometallurgy, or mine digitization separately, but to connect those literatures and argue that DRP can serve as a unifying technical layer between exploration, process design, recycling, and porous-material innovation. By redirecting this expertise from hydrocarbons to hard rocks and combining multiple imaging modalities to characterize both ores and engineered porous materials, the region can improve how it explores, processes, and recycles mineral resources while strengthening supply-chain resilience for the energy transition.

\begin{thebibliography}{99}

\bibitem{drp_blunt_2013}
Blunt, M.J., Bijeljic, B., Dong, H., Gharbi, O., Iglauer, S., Mostaghimi, P., Paluszny, A., \& Pentland, C., ``Pore-scale imaging and modelling,'' \textit{Advances in Water Resources}, 51, 197--216, 2013. DOI: \url{https://doi.org/10.1016/j.advwatres.2012.03.003}.

\bibitem{lishchuk_predictive_2019}
Lishchuk, V., Lund, C., \& Ghorbani, Y., ``Evaluation and comparison of different machine-learning methods to integrate sparse process data into a spatial model in geometallurgy,'' \textit{Minerals Engineering}, 138, 144--153, 2019. DOI: \url{https://doi.org/10.1016/j.mineng.2019.01.032}.

\bibitem{ai_mineral_prospectivity_2011}
Zuo, R., \& Carranza, E.J.M., ``Support vector machine: A tool for mapping mineral prospectivity,'' \textit{Computers \& Geosciences}, 37(12), 1967--1975, 2011. DOI: \url{https://doi.org/10.1016/j.cageo.2010.09.014}.

\bibitem{norgate_comminution_2010}
Norgate, T., \& Haque, N., ``Energy and greenhouse gas impacts of mining and mineral processing operations,'' \textit{Journal of Cleaner Production}, 18(3), 266--274, 2010. DOI: \url{https://doi.org/10.1016/j.jclepro.2009.09.020}.

\bibitem{kyle_processing_2015}
Kyle, J.R., \& Ketcham, R.A., ``Application of high resolution X-ray computed tomography to mineral deposit origin, evaluation, and processing,'' \textit{Ore Geology Reviews}, 65, 821--839, 2015. DOI: \url{https://doi.org/10.1016/j.oregeorev.2014.09.034}.

\bibitem{mri_heap_leach_2015}
Fagan-Endres, M.A., Harrison, S.T.L., Johns, M.L., \& Sederman, A.J., ``Magnetic resonance imaging characterisation of the influence of flowrate on liquid distribution in drip irrigated heap leaching,'' \textit{Hydrometallurgy}, 157, 116--124, 2015. DOI: \url{https://doi.org/10.1016/j.hydromet.2015.10.003}.

\bibitem{mofs_lithium_2026}
Cao, X., Gao, H., \& Zeng, Y., ``Metal-organic frameworks for extraction of radionuclide/noble metals/lithium,'' in \textit{Metal Organic Frameworks for Energy and Environmental Applications}, 13--33, 2026. DOI: \url{https://doi.org/10.1016/B978-0-443-40377-4.00002-1}.

\bibitem{araya_tailings_2020}
Araya, N., Kraslawski, A., \& Cisternas, L.A., ``Towards mine tailings valorization: Recovery of critical materials from Chilean mine tailings,'' \textit{Journal of Cleaner Production}, 263, 121555, 2020. DOI: \url{https://doi.org/10.1016/j.jclepro.2020.121555}.

\bibitem{battery_blackmass_2021}
Vanderbruggen, A., Gugala, E., Blannin, R., Bachmann, K., Serna-Guerrero, R., \& Rudolph, M., ``Automated mineralogy as a novel approach for the compositional and textural characterization of spent lithium-ion batteries,'' \textit{Minerals Engineering}, 169, 106924, 2021. DOI: \url{https://doi.org/10.1016/j.mineng.2021.106924}.

\bibitem{buyse_correlative_2023}
Buyse, F., Dewaele, S., Boone, M.N., \& Cnudde, V., ``Combining Automated Mineralogy with X-ray Computed Tomography for Internal Characterization of Ore Samples at the Microscopic Scale,'' \textit{Natural Resources Research}, 32(2), 461--478, 2023. DOI: \url{https://doi.org/10.1007/s11053-023-10161-z}.

\bibitem{graczyk_dl_porosity_2020}
Graczyk, K.M., \& Matyka, M., ``Predicting porosity, permeability, and tortuosity of porous media from images by deep learning,'' \textit{Scientific Reports}, 10, 21459, 2020. DOI: \url{https://doi.org/10.1038/s41598-020-78415-x}.

\bibitem{karniadakis_piml_2021}
Karniadakis, G.E., Kevrekidis, I.G., Lu, L., Perdikaris, P., Wang, S., \& Yang, L., ``Physics-informed machine learning,'' \textit{Nature Reviews Physics}, 3, 422--440, 2021. DOI: \url{https://doi.org/10.1038/s42254-021-00314-5}.

\bibitem{raissi_pinns_2019}
Raissi, M., Perdikaris, P., \& Karniadakis, G.E., ``Physics-informed neural networks: A deep learning framework for solving forward and inverse problems involving nonlinear partial differential equations,'' \textit{Journal of Computational Physics}, 378, 686--707, 2019. DOI: \url{https://doi.org/10.1016/j.jcp.2018.10.045}.

\bibitem{ml_upscaling_drp_2024}
You, J., \& Lee, K.J., ``Upscaling reactive transport models from pore-scale to continuum-scale using deep learning method,'' \textit{Geoenergy Science and Engineering}, 242, 212850, 2024. DOI: \url{https://doi.org/10.1016/j.geoen.2024.212850}.

\bibitem{matter_carbfix_2009}
Matter, J.M., Broecker, W.S., Stute, M., Gislason, S.R., Oelkers, E.H., Stefansson, A., \textit{et al.}, ``Permanent Carbon Dioxide Storage into Basalt: The CarbFix Pilot Project, Iceland,'' \textit{Energy Procedia}, 1(1), 3641--3646, 2009. DOI: \url{https://doi.org/10.1016/j.egypro.2009.02.160}.

\bibitem{kelemen_peridotite_2020}
Kelemen, P.B., McQueen, N., Wilcox, J., Renforth, P., Dipple, G., \& Vankeuren, A.P., ``Engineered carbon mineralization in ultramafic rocks for CO2 removal from air: Review and new insights,'' \textit{Chemical Geology}, 550, 119628, 2020. DOI: \url{https://doi.org/10.1016/j.chemgeo.2020.119628}.

\bibitem{menke_hydrogen_2024}
Jangda, Z., Menke, H., Busch, A., Geiger, S., Bultreys, T., \& Singh, K., ``Subsurface hydrogen storage controlled by small-scale rock heterogeneities,'' \textit{International Journal of Hydrogen Energy}, 72, 467--482, 2024. DOI: \url{https://doi.org/10.1016/j.ijhydene.2024.02.233}.

\bibitem{bartels_fast_xray_2017}
Bartels, W.-B., Lin, Q., de Winter, D.A.M., Mahani, H., Berg, S., Cnudde, V., \& Ajo-Franklin, J.B., ``Fast X-Ray Micro-CT Study of the Impact of Brine Salinity on the Pore-Scale Fluid Distribution During Waterflooding,'' \textit{Petrophysics}, 58(1), 36--47, 2017.

\bibitem{farber_science_2025}
Farber, E.M., Seraphim, N.M., Tamakuwala, K., Stein, A., Rücker, M., \& Eisenberg, D., ``Porous materials: The next frontier in energy technologies,'' \textit{Science}, 387(6730), eadn9391, 2025. DOI: \url{https://doi.org/10.1126/science.adn9391}.

\bibitem{digitalrocks_portal_2023}
Prodanović, M., Esteva, M., McClure, J., Chang, B.C., Santos, M., Radhakrishnan, S., \textit{et al.}, ``Digital Rocks Portal (Digital Porous Media): Connecting data, simulation and community,'' \textit{E3S Web of Conferences}, 367, 01010, 2023. DOI: \url{https://doi.org/10.1051/e3sconf/202336701010}.

\bibitem{calvo_copper_2016}
Calvo, G., Mudd, G., Valero, A., \& Valero, A., ``Decreasing Ore Grades in Global Metallic Mining: A Theoretical Issue or a Global Reality?,'' \textit{Resources}, 5(4), 36, 2016. DOI: \url{https://doi.org/10.3390/resources5040036}.

\bibitem{iea_outlook_2025}
International Energy Agency (IEA), \textit{Global Critical Minerals Outlook 2025}, IEA, Paris, 2025. \url{https://www.iea.org/reports/global-critical-minerals-outlook-2025}.

\bibitem{eu_crma_2024}
European Union, ``Regulation (EU) 2024/1252 of the European Parliament and of the Council of 11 April 2024 establishing a framework for ensuring a secure and sustainable supply of critical raw materials,'' \textit{Official Journal of the European Union}, 2024. \url{https://eur-lex.europa.eu/eli/reg/2024/1252/oj/eng}.

\bibitem{uk_strategy_2022}
HM Government, \textit{Resilience for the Future: The UK's Critical Minerals Strategy}, London, 2022. \url{https://www.gov.uk/government/publications/uk-critical-mineral-strategy}.

\bibitem{doe_arpae}
Advanced Research Projects Agency-Energy (ARPA-E), \textit{ARPA-E FY 2023 Annual Report to Congress}, U.S. Department of Energy, August 2025. \url{https://arpa-e.energy.gov/sites/default/files/2025-09/ARPA-E%20FY%202023%20Annual%20Report.pdf}.

\bibitem{doe_mine_future_2025}
U.S. Department of Energy, ``Funding Notice: Mine of the Future -- Proving Ground Initiative,'' Hydrocarbons and Geothermal Energy Office, November 14, 2025. \url{https://www.energy.gov/hgeo/funding-notice-mine-future-proving-ground-initiative}.

\bibitem{netl_digitalrocks}
National Energy Technology Laboratory (NETL), \textit{Digital Rocks Portal}, U.S. Department of Energy. \url{https://digitalporousmedia.org/digital-rocks-portal-tombstone/}. Accessed March 24, 2026.

\bibitem{diamond_tomography_2026}
Diamond Light Source, ``Tomography,'' \textit{Imaging Techniques}, \url{https://www.diamond.ac.uk/Instruments/Techniques/Imaging/Tomography.html}. Accessed March 24, 2026.

\bibitem{esrf_bm05_2026}
European Synchrotron Radiation Facility (ESRF), ``BM05 - Industrial Imaging beamline,'' \url{https://www.esrf.fr/fr/home/UsersAndScience/Experiments/BM05.html}. Accessed March 24, 2026.

\bibitem{psi_tomcat_2025}
Paul Scherrer Institute (PSI), ``Beamline Information | BL: SLS/TOMCAT,'' \url{https://www.psi.ch/de/sls/tomcat/beamlines}. Accessed March 24, 2026.

\bibitem{yang_ai_exploration_2024}
Yang, L., Zuo, R., \& Kreuzer, O.P., ``Artificial intelligence for mineral exploration: A review and perspectives on future directions from data science,'' \textit{Earth-Science Reviews}, 255, 104941, 2024. DOI: \url{https://doi.org/10.1016/j.earscirev.2024.104941}.

\bibitem{sarker_tailings_2022}
Sarker, P.K., Haque, N., Bhuiyan, M.A.H., Bruckard, W., \& Pramanik, B.K., ``Recovery of strategically important critical minerals from mine tailings,'' \textit{Journal of Environmental Chemical Engineering}, 10(3), 107622, 2022. DOI: \url{https://doi.org/10.1016/j.jece.2022.107622}.

\bibitem{peukert_sorting_2022}
Peukert, W., Xu, Z., \& Dowd, P.A., ``A Review of Sensor-Based Sorting in Mineral Processing: The Potential Benefits of Sensor Fusion,'' \textit{Minerals}, 12(11), 1364, 2022. DOI: \url{https://doi.org/10.3390/min12111364}.

\bibitem{mojid_dle_2024}
Mojid, R.A., Lee, M.-G., \& You, J.-K., ``A review on advances in direct lithium extraction from continental brines: Ion-sieve adsorption and electrochemical methods for varied Mg/Li ratios,'' \textit{Sustainable Materials and Technologies}, 41, e00923, 2024. DOI: \url{https://doi.org/10.1016/j.susmat.2024.e00923}.

\bibitem{baird_dle_2024}
Baird, E.J., Kingsbury, R.S., Wang, X., Gang, O., Zhu, Y., \& Whittaker, M.L., ``Microporous Polymer Sorbents for Direct Lithium Extraction,'' \textit{ACS Energy Letters}, 9(9), 4229--4237, 2024. DOI: \url{https://doi.org/10.1021/acsenergylett.4c01590}.

\bibitem{koch_geomet_2019}
Koch, P.H., Lund, C., \& Rosenkranz, J., ``Automated drill core mineralogical characterization method for texture classification and modal mineralogy estimation for geometallurgy,'' \textit{Minerals Engineering}, 136, 99--110, 2019. DOI: \url{https://doi.org/10.1016/j.mineng.2019.03.008}.

\bibitem{gu_automated_mineralogy_2014}
Gu, Y., Schouwstra, R.P., \& Rule, C., ``The value of automated mineralogy,'' \textit{Minerals Engineering}, 58, 100--103, 2014. DOI: \url{https://doi.org/10.1016/j.mineng.2014.01.020}.

\bibitem{reich_critical_2025}
Reich, M., \& Simon, A.C., ``Critical Minerals,'' \textit{Annual Review of Earth and Planetary Sciences}, 53, 149--177, 2025. DOI: \url{https://doi.org/10.1146/annurev-earth-040523-023316}.

\bibitem{lishchuk_integrated_2020}
Lishchuk, V., Koch, P.H., Ghorbani, Y., \& Butcher, A.R., ``Towards integrated geometallurgical approach: Critical review of current practices and future trends,'' \textit{Minerals Engineering}, 145, 106072, 2020. DOI: \url{https://doi.org/10.1016/j.mineng.2019.106072}.

\bibitem{becker_applied_2023}
Becker, M., ``The contribution of applied mineralogy to sustainability in the mine life cycle,'' \textit{Minerals Engineering}, 202, 108121, 2023. DOI: \url{https://doi.org/10.1016/j.mineng.2023.108121}.

\bibitem{nwaila_minewaste_2021}
Nwaila, G.T., Ghorbani, Y., Zhang, S.-E., Tolmay, V., Rose, D.H., \& Nwaila, N.G., ``Valorisation of mine waste - Part II: Resource evaluation for consolidated and mineralised mine waste using the Central African Copperbelt as an example,'' \textit{Journal of Environmental Management}, 288, 113553, 2021. DOI: \url{https://doi.org/10.1016/j.jenvman.2021.113553}.

\bibitem{guntoro_liberation_2021}
Guntoro, P.I., Ghorbani, Y., Parian, M., Butcher, A.R., Kuva, J., \& Rosenkranz, J., ``Development and experimental validation of a texture-based 3D liberation model,'' \textit{Minerals Engineering}, 173, 106828, 2021. DOI: \url{https://doi.org/10.1016/j.mineng.2021.106828}.

\end{thebibliography}

\end{document}
